\documentclass[12pt]{article}
\usepackage[utf8]{inputenc}
\usepackage[T1]{fontenc}
\usepackage[italian]{babel}

\title{Peer-Review : Network}
\author{Riva Filippo, Vinci Francesco, Sulfaro Antonio, Sonzini Anna\\Gruppo 16}

\begin{document}

\maketitle

Valutazione dello strato di rete del gruppo 06.

\section{Lati Positivi}

I punti di forza di questa progettazione :
\begin{enumerate}
    \item La semplicità della modellazione. Nonostante il ridotto numero di classi e di messaggi questo approccio è in grado di soddisfare tutte le specifiche e le funzionalità aggiuntive scelte, con un limitato numero di metodi e senza complicare eccessivamente la logica all'interno e all'esterno dello strato di rete.
    \item La leggibilità dell'UML. 
    \item Un limitato numero di messaggi. 
\end{enumerate}

\section{Lati negativi}

I punti deboli che abbiamo riscontrato nella progettazione sono : 
\begin{enumerate}
    \item Un UML non sufficientemente dettagliato alla comprensione dello scambio di messaggi, in quanto non specifica cosa essi possano contenere e come il ricevente sia in grado di interpretarli. 
    \item Manca un diagramma di flusso per comprendere meglio come viene sfruttato il pattern observer (che nell'UML è ovviamente difficile da rappresentare in assenza di classi che lo implementano) anche se in ogni caso è stato spiegato nel commento all'UML.
    \item Invio dell'oggetto Client ad ogni messaggio del Client verso il Server, appesantendo inutilmente i messaggi quando basterebbe mandare il nome del giocatore.
\end{enumerate}

\section{Confronto tra le architetture}

Il problema di progettazione da parte del nostro gruppo è stato affrontato con un approccio simile. 
Le principali differenze sono :
\begin{itemize}
    \item La gestione delle connessioni, in quanto nel nostro caso non vi è un unico oggetto server a cui mandare i messaggi ma ogni client parla al proprio oggetto ServerConnectionHandler lato server che ne gestisce i messaggi.  
    \item Il sistema di ping utilizzato dal client per notificare il server della propria connessione, utile alla funzionalità avanzata della resilienza alle disconnessioni. 
\end{itemize}


Dei possibili miglioramenti del nostro strato di rete potrebbero essere :

\begin{itemize}
    \item Una riduzione del numero di messaggi e delle classi rendendone più facile la comprensione ed eventuali modifiche future.
    \item L'implementazione di un sistema di rilevazione delle disconnessioni tramite ping o simili. 
\end{itemize}

\end{document}
